%% notes-tex/031-unified-representation/sections/3b-joinability.tex

\section{Can these datasets be joined at all\secstatus{todo}}
\label{sec:joinability}

Pooling is an assertion about the LABEL, not only about the input. Two records that measure
the same gene under the same compound have to mean the same thing before one loss can score
them both. This section tests that assertion on the served records rather than assuming it,
and the test changes what the pooled target has to look like.

A cell here is one (gene, compound) pair. A record enters only if it doses exactly one
compound carrying an InChIKey and perturbs exactly one gene, so nothing is averaged across
combinations; Vanacloig's three host deletions are stripped first, leaving the queried gene.

%% SOURCE: experiments/031-env-chemgen-inhibitor-tolerance/scripts/dataset_joinability.py ->
%% results/dataset_distributions.csv and results/cross_dataset_pair_overlap.csv, rendered by
%% notes_tex_representation_tables.py.
\begin{table}[htbp]
\centering
\footnotesize
\caption{What each dataset measures and reports, and the sign its source declares for a sick
strain. The standard deviation and skew are of the served response over its cells. The last
column is the factor the orientation step applies.}
\label{tab:unification}
%% GENERATED by experiments/031-env-chemgen-inhibitor-tolerance/scripts/notes_tex_representation_tables.py
%% SOURCE: experiments/031-env-chemgen-inhibitor-tolerance/results/dataset_distributions.csv, genotype_overlap.csv and dose_axis_summary.csv from dataset_joinability.py, exploratory_overlap_panels.py and environment_axis_coverage.py
%% Do not edit by hand.
\begin{tabular}{lllllrrll}
\toprule
\textbf{dataset} & \textbf{genes x compounds} & \textbf{cells} & \textbf{dose basis} & \textbf{reports} & \textbf{sd} & \textbf{skew} & \textbf{sick} & \textbf{orient} \\
\midrule
Vanacloig 2022 & 3,598 x 41 & 143,218 & IC30, fixed & log2 ratio & 0.70 & -2.6 & negative & -1 \\
Hillenmeyer HOM & 4,655 x 114 & 483,279 & fixed & z score & 2.35 & +5.6 & positive & +1 \\
Hillenmeyer HET & 5,810 x 290 & 1,651,248 & fixed & log2 ratio & 0.38 & +1.2 & positive & +1 \\
Hoepfner 2014 & 5,839 x 148 & 861,345 & IC30 & sensitivity score & 1.29 & -4.4 & negative & -1 \\
Wildenhain 2015 & 242 x 5,170 & 428,189 & none & z score & 7.50 & -4.6 & negative & -1 \\
\bottomrule
\end{tabular}

\end{table}

\notesec{The scales differ by a factor of twenty}
The served standard deviations run from 0.38 for Hillenmeyer HET to 7.50 for Wildenhain
(Table~\ref{tab:unification}, Figure~\ref{fig:joinability}a and b). A pooled squared-error
loss weights a record by the square of its residual in native units, so without a per-source
scale the objective is dominated by Wildenhain and by the tail of Hillenmeyer HOM, whatever
the record counts are. Standardizing each source by its own training-split statistics is
therefore not a cosmetic step, and the statistics come from the training split alone so that
no held-out value informs the scale.

\notesec{The polarity is not shared, and this is the finding that matters}
A more negative number does not mean a sicker strain in every dataset. Each loader records
its source's own definition verbatim, and those definitions split the five into two groups.
Vanacloig stores $\log_2(\text{inhibitor}/\text{control})$, so depletion is negative. Hoepfner
stores an adjusted MADL sensitivity score whose definition reads ``negative = hypersensitive,
positive = resistant''. Wildenhain stores a growth-inhibition z score where negative is
inhibition. Both Hillenmeyer arms store a fitness-defect statistic whose definition reads
``positive = fitness defect''.

Two independent checks confirm the split rather than leaving it as a reading of the
documentation. First, the heavy tail of each distribution falls on the side its source names:
the three negative-is-sick datasets have skew $-2.6$, $-4.4$ and $-4.6$, and the two
positive-is-sick arms have skew $+5.6$ and $+1.2$ (Figure~\ref{fig:joinability}c). Second,
the declared polarities predict the SIGN of every pairwise correlation between datasets, and
they get all ten right (Table~\ref{tab:crossoverlap}). Under a null in which the declared
polarity says nothing about the data, ten correct signs out of ten has probability $0.002$.

The consequence is concrete and large. As served, Hillenmeyer HOM and Hoepfner correlate at
$-0.198$ over 100{,}276 shared cells, and Hillenmeyer HET and Hoepfner at $-0.118$ over
170{,}837. Pooling without orienting trains a model on contradictory labels at that scale,
and the contradiction is worst exactly where the overlap is largest.

%% SOURCE: experiments/031-env-chemgen-inhibitor-tolerance/scripts/dataset_joinability.py ->
%% results/cross_dataset_pair_overlap.csv, rendered by notes_tex_representation_tables.py.
\begin{table}[htbp]
\centering
\footnotesize
\caption{Every dataset pair: the genes, compounds and (gene, compound) cells both measure,
the Spearman correlation on those cells as served, the sign predicted from the two declared
polarities, and the correlation after orienting. The predicted sign is correct in all ten
pairs.}
\label{tab:crossoverlap}
%% GENERATED by experiments/031-env-chemgen-inhibitor-tolerance/scripts/notes_tex_representation_tables.py
%% SOURCE: experiments/031-env-chemgen-inhibitor-tolerance/results/cross_dataset_pair_overlap.csv from dataset_joinability.py
%% Do not edit by hand.
\begin{tabular}{lrrrrrr}
\toprule
\textbf{pair} & \textbf{genes} & \textbf{compounds} & \textbf{cells} & \textbf{$\rho$ as served} & \textbf{sign predicted} & \textbf{$\rho$ oriented} \\
\midrule
Hillenmeyer HOM / Hillenmeyer HET & 4,625 & 99 & 407,072 & +0.117 & +1 & 0.117 \\
Hillenmeyer HET / Hoepfner 2014 & 5,802 & 30 & 170,837 & -0.118 & -1 & 0.118 \\
Hillenmeyer HOM / Hoepfner 2014 & 4,651 & 23 & 100,276 & -0.198 & -1 & 0.198 \\
Hillenmeyer HET / Wildenhain 2015 & 238 & 123 & 17,218 & -0.018 & -1 & 0.018 \\
Vanacloig 2022 / Hoepfner 2014 & 3,592 & 2 & 6,977 & +0.013 & +1 & 0.013 \\
Vanacloig 2022 / Hillenmeyer HOM & 3,544 & 2 & 6,854 & -0.075 & -1 & 0.075 \\
Vanacloig 2022 / Hillenmeyer HET & 3,571 & 2 & 6,852 & -0.023 & -1 & 0.023 \\
Hoepfner 2014 / Wildenhain 2015 & 237 & 34 & 4,787 & +0.061 & +1 & 0.061 \\
Hillenmeyer HOM / Wildenhain 2015 & 199 & 30 & 3,292 & -0.001 & -1 & 0.001 \\
Vanacloig 2022 / Wildenhain 2015 & 153 & 10 & 761 & +0.078 & +1 & 0.078 \\
\bottomrule
\end{tabular}

\end{table}

%% SOURCE: experiments/031-env-chemgen-inhibitor-tolerance/scripts/dataset_joinability.py ->
%% dataset_joinability.svg, converted by `make plots`.
\tcfig{figures/dataset_joinability.pdf}{Whether the five datasets can share one target.
\textbf{(a)} Response density per dataset on a log scale, which is where the phenotype lives.
\textbf{(b)} Standard deviation of the served response. \textbf{(c)} Skew beside the sign each
source declares for a sick strain; hatching marks the two datasets whose sources use the
opposite convention. \textbf{(d)} Shared (gene, compound) cells per pair, log color scale.
\textbf{(e)} Spearman on those shared cells after orienting. \textbf{(f)} The same pairs as
served against oriented, point area by shared cells; red points are pairs whose sources
disagree about the sign.}{fig:joinability}

\notesec{Agreement after orienting is real, positive, and small}
Once oriented, every pair correlates positively, but weakly: the median over the ten pairs is
$0.068$ and the largest is $0.198$, for Hillenmeyer HOM against Hoepfner over 100{,}276 cells
(Table~\ref{tab:crossoverlap}, Figure~\ref{fig:joinability}e). The two Hillenmeyer arms, which
share a laboratory, a medium and a readout family, reach $0.117$ over 407{,}072 cells. Against
the application target the numbers are smaller still, because Vanacloig shares only two
compounds with each of the Hillenmeyer arms and with Hoepfner, and ten with Wildenhain.

That is the honest shape of the case for joining. These datasets agree about which genes are
generally fragile far more than about how a particular gene answers a particular compound. So
a shared SCALE is justified and a shared RESPONSE is not: pooling should be expected to buy a
gene prior and a wider compound space, not more labels for the same task. Section~\ref{sec:model}
turns that into a decoder that gives each source its own token, and Section~\ref{sec:established}
already measured the alternative, which is that within one dataset, one medium and one dose
basis, the same molecule features reach $0.31$ where they reach at most $0.09$ across
datasets.

%% SOURCE: notes/experiments.031-env-chemgen-inhibitor-tolerance.mermaid.dataset-unification.md,
%% rendered by notes/assets/publish/scripts/mermaid_pdf.sh through `make diagrams`.
\tcfig{figures/dataset_unification.pdf}{One row per kept dataset: what the record is, what the
source reports, and the common form both are mapped into. The input tuple and the scalar
target are identical for every row. The only per-source objects are the declared sick-sign
$s_k$, the training-split scale $\mathrm{sd}_k$ and the source token $t$. The per-dataset
counts are in Table~\ref{tab:unification}.}{fig:unification}

\notesec{Where the datasets overlap on every other axis}
Figure~\ref{fig:overlap} puts the same question to the genotype, chemistry and dose axes.
Genes overlap heavily and compounds barely: 98 to 100 percent of Vanacloig's genes appear in
each of the three genome-wide partners, while at most 24 percent of its compounds appear in
any one of them. Wildenhain is the exception in both directions, covering 4 percent of the
other gene universes and contributing by far the most chemistry. The chemical panels show why
the compound gain is uneven: the partner libraries are drug-like at 20 to 40 heavy atoms, and
the hydrolysate panel sits at 3 to 15, in a region where a substructure fingerprint is nearly
empty. The dose panel repeats Section~\ref{sec:dose} visually, with Vanacloig carrying no
convertible concentration at all.

%% SOURCE: experiments/031-env-chemgen-inhibitor-tolerance/scripts/exploratory_overlap_panels.py ->
%% exploratory_overlap.svg, converted by `make plots`.
\tcfig{figures/exploratory_overlap.pdf}{Overlap on every axis a pooled model sees.
\textbf{(a)} Share of the row dataset's genes that the column dataset also measures, with each
dataset's gene count. \textbf{(b)} The same for compounds. \textbf{(c)} Molecular weight against
calculated logP for every dosed compound, the Vanacloig panel outlined. \textbf{(d)} Heavy
atom count; the shaded band is where a substructure fingerprint carries almost no bits.
\textbf{(e)} Each Vanacloig compound's best non-exact Tanimoto in the partner union against its
size. \textbf{(f)} The molar range each dataset spans, 1st to 99th percentile, with the share
of records stating a convertible dose. \textbf{(g)} Distinct compounds against distinct genes,
point area by measured cells. \textbf{(h)} Perturbation class by dataset.}{fig:overlap}

